\chapter{Fragmented Apertures Applied to Other Electromagnetic Structures}

\section{Introduction}

The preceding chapters have focused exclusively on the application of the fragmented aperture design methodology to antennas---radiating structures that convert guided electromagnetic energy into free-space radiation (or the reverse).  However, the core idea behind the fragmented aperture approach---partitioning a conducting surface into sub-wavelength fragments and using computational optimization to determine the optimal pattern---is not limited to antenna design.  The same methodology can be applied to any electromagnetic structure whose performance depends on the spatial distribution of conducting material.

This chapter surveys several non-antenna applications of the fragmented aperture concept.  These include broadband phase surfaces, RF power combining networks, fragmented absorbers, fragmented waveguide components, and metamaterials.  In each case, the same design machinery described in the earlier chapters---a fragmented conducting surface, a genetic algorithm (or other optimizer), and a full-wave electromagnetic simulation---is applied to a structure whose function is fundamentally different from that of an antenna.

\section{Fragmented Broadband Phase Surfaces}
\label{sec:phaseSurface}

An early and important application of the fragmented aperture concept was the design of surfaces with broadband transmission phase properties.  Just as the fragment pattern on a radiating aperture can be optimized to achieve desired antenna characteristics, the fragment pattern on a standalone surface can be optimized to achieve a specific phase response as electromagnetic waves pass through it.

A conventional solid conducting surface reflects all incident energy and imparts a frequency-dependent phase shift upon reflection.  Similarly, transmission through a partially transparent surface depends on the size, shape, and arrangement of the conducting regions.  For a solid conductor, the transmission phase varies rapidly with frequency.  By fragmenting the conducting surface into an optimized pattern of conducting and non-conducting fragments, it is possible to design a surface that provides a more uniform phase response over a wider bandwidth.

Figure~\ref{fig:OES1} shows the transmission phase through a fragmented surface (curve 18) compared with a reference solid conducting surface (curve 19), from the original patent \cite{MaloneyFragPatent}.  The fragmented surface maintains a substantially flatter phase response across the normalized frequency range, demonstrating that the fragmented pattern can be used to engineer broadband phase behavior.

\begin{figure}
\includegraphics[angle=0.5,trim={0 1100bp 0 0},clip,width=\linewidth]{/Users/jim.maloney/Book/images/ch2/origPatentFig2.png}
\caption[Broadband transmission phase of fragmented surface]{Transmission phase comparison demonstrating the broadband properties of a fragmented surface.  The fragmented surface (curve 18) maintains a flatter phase response than the reference solid conductor (curve 19) over a wide normalized frequency range \cite{MaloneyFragPatent}.}
\label{fig:OES1}
\end{figure}

This result was significant because it demonstrated that the fragmented aperture concept extends beyond antenna design to the engineering of electromagnetic surfaces with tailored transmission and reflection properties.  The broadband phase behavior of fragmented surfaces subsequently informed the development of the broadband screen backplanes described in Chapter~6, which use fragmented conducting layers in combination with resistive card (r-card) layers to achieve ultra-wideband ground plane performance for phased array antennas.

The concept of designing surfaces with optimized phase properties is closely related to the broader field of frequency-selective surfaces (FSS) and metasurfaces.  As described in Chapter~9, Saad-Falcon et al.\ \cite{SaadFalcon2024LevelSet} applied the Periodic Level Set Function parameterization to design broadband bandpass and narrowband notch filter metasurfaces, achieving excellent agreement between simulation and measurement.  Howard et al.\ \cite{Howard2022LossTangent} applied similar techniques to design resonant periodic surfaces for dielectric materials characterization, where the resonance magnitude depends linearly on the loss tangent of an adjacent dielectric sample.  These recent developments demonstrate that the fragmented aperture approach to surface design has matured from the original broadband phase concept into a versatile tool for engineering a wide range of electromagnetic surface behaviors.


\section{Fragmented Reflectarrays}
\label{sec:reflectarrays}

Reflectarray antennas combine the high gain of parabolic reflectors with the low profile of planar arrays: an array of reflecting elements on a flat surface is illuminated by a feed horn, and each element imparts a specific reflection phase to steer and focus the beam.  Conventional reflectarray elements (patches, rings, stubs) provide phase control through a single geometric parameter, but the limited phase range and narrow bandwidth of these canonical shapes constrain overall performance.  Ethier et al.\ \cite{Ethier2014Reflectarray} demonstrated that fragmenting each reflectarray element into a sub-wavelength fragment grid and optimizing the fragment pattern with a genetic algorithm yields a powerful alternative approach.

\begin{figure}
\begin{center}
\includegraphics[angle=0,width=\linewidth]{/Users/jim.maloney/Book/images/ch10/ethier2014/Ethier2014_Fig4_SimilarityElements.png}
\caption[Similarity-optimized fragmented reflectarray elements]{Set of 31 similarity-optimized fragmented reflectarray elements, each with its reflection phase indicated.  Each element is a 19$\times$19 fragment grid ($\lambda/6$ cell at 12~GHz) with quadrantal symmetry.  The GA optimizes each element's fragment pattern to achieve the target phase while maximizing geometric similarity to neighboring elements in the array \cite{Ethier2014Reflectarray}.}
\label{fig:OES_reflectarray1}
\end{center}
\end{figure}

Each element in the reflectarray occupies a 4~mm cell ($\sim\!\lambda/6$ at 12~GHz) subdivided into a 19$\times$19 grid of 0.2~mm fragments with quadrantal symmetry.  The key innovation is \emph{similarity optimization}: the GA shapes each element to provide the required reflection phase while maximizing geometric similarity to its neighbors, measured by an XNOR-based metric.  This ensures that the aperiodic reflectarray more closely resembles the infinite periodic assumption used in element design, reducing phase errors at Fresnel zone boundaries (Figure~\ref{fig:OES_reflectarray1}).

\begin{figure}
\begin{center}
\includegraphics[angle=0,width=\linewidth]{/Users/jim.maloney/Book/images/ch10/ethier2014/Ethier2014_Fig12_FabricatedArray.png}
\caption[Fabricated fragmented reflectarray aperture]{Fabricated similarity-optimized fragmented reflectarray (Design~D) with 9{,}801 elements on a 420$\times$420~mm aperture, indicating the 4~mm cell dimension.  Each element is a unique fragmented pattern optimized for its required reflection phase \cite{Ethier2014Reflectarray}.}
\label{fig:OES_reflectarray2}
\end{center}
\end{figure}

A 99$\times$99 element reflectarray (9{,}801 elements) was fabricated on FR4 substrate backed by 4~mm Rohacell foam over an aluminum ground plane (Figure~\ref{fig:OES_reflectarray2}).  The similarity-optimized design achieved 33.1~dBi gain at 11.2~GHz with 12.3\% bandwidth (1~dB), first sidelobe level of $-20.6$~dB, and an aperture phase efficiency of approximately 97\%---far exceeding the 40--55\% typical of conventional reflectarrays.  The dissimilarity-optimized counterpart, using the same fragment parameterization but without the similarity constraint, achieved 0.9~dB lower gain and 3.4~dB worse sidelobe performance, confirming the importance of inter-element similarity.

The authors also demonstrated a ``mosaic reflectarray'' concept: by controlling both the reflection phase and the fragment density (metal fill fraction) of each element, the reflectarray surface can simultaneously display a grayscale optical image while functioning as a high-gain antenna.  A mosaic reflectarray displaying the University of Ottawa logo achieved 32.4~dBi gain---only 0.7~dB below the similarity-optimized design---demonstrating that fragmented reflectarray elements offer sufficient design freedom to satisfy both electromagnetic and visual constraints simultaneously.


\section{Fragmented RF Power Combiners}
\label{sec:combiners}

An important practical challenge in phased array antenna systems is the design of the power combining network that combines the signals from multiple array elements into a single output port (for receive) or distributes a single input signal to multiple elements (for transmit).  Conventional power combiners use commercial off-the-shelf (COTS) components---Wilkinson dividers, hybrid couplers, or resistive dividers---connected by microstrip or stripline transmission lines.  These COTS networks can be bulky, expensive, and lossy, particularly when they must operate over wide bandwidths.

Landgren et al.\ at GTRI recognized that the fragmented aperture optimization process could be applied to the design of power combiner networks, not just radiating surfaces.  The key insight is that a power combiner is simply another electromagnetic structure whose performance---insertion loss, reflection, isolation, and bandwidth---depends on the spatial distribution of conducting material on a dielectric substrate.  By treating the combiner as a fragmented conducting pattern and optimizing it with a genetic algorithm and FDTD simulation, the same design machinery used for fragmented aperture antennas can produce optimized combiner networks.

\subsection{Fragmented Combiner Concept}

Landgren et al.\ \cite{Landgren2018Combiner} demonstrated this concept by designing a fragmented combiner network for a 4$\times$4 phased array antenna intended for fire-control radar telemetry applications, covering the $L$-, $S$-, and $C$-bands (1.4--6.7~GHz, a 4.8:1 bandwidth).  The combiner network replaces the conventional Wilkinson divider tree with a single printed circuit board layer containing an optimized fragmented conducting pattern.

The fragmented combiner design required two separate optimized components: an interfacial combiner that interfaces the unbalanced antenna feed structure with a microstrip transmission line, and an internal combiner that provides the actual power division.  A collection of these two logical units was arranged to form a 16:1 combiner for the 4$\times$4 array.  The design process used a genetic algorithm coupled with FDTD simulation, with the fitness function targeting minimum insertion loss, low reflection, and high isolation across the 1.4--6.7~GHz band.

Figure~\ref{fig:OES2} shows the key innovation: the conventional COTS Wilkinson combiner network (left) uses discrete components and transmission line segments, while the fragmented combiner network (right) replaces all of these with a single optimized conducting pattern on a PCB substrate.  Both are designed for a 4$\times$4 array with a 16:1 combining ratio and the same physical footprint.

\begin{figure}
\begin{center}
\includegraphics[angle=0,width=\linewidth]{/Users/jim.maloney/Book/images/ch10/Landgren2018_Fig1_cropped.png}
\caption[Wilkinson vs.\ fragmented combiner network]{Comparison of a conventional COTS Wilkinson combiner network (left) and a fragmented combiner network (right), both designed for a 4$\times$4 array with a 16:1 combining ratio.  The fragmented combiner replaces discrete components with a single optimized conducting pattern \cite{Landgren2018Combiner}.}
\label{fig:OES2}
\end{center}
\end{figure}

\subsection{Measured Results}

The fragmented combiner was fabricated on Rogers 4350B substrate and experimentally characterized.  The measured insertion loss of the fragmented combiner was lower than that of a conventional COTS Wilkinson combiner across the entire operating band.  The ideal insertion loss for a 16:1 combiner is $-12$~dB ($-3$~dB per 2:1 stage); the fragmented combiner was within 2.5~dB of this ideal, compared to larger losses for the COTS version.

Figure~\ref{fig:OES3} shows photographs of the fabricated combiner boards and the measured insertion loss comparison.  The COTS combiner (a) has Wilkinson divider chips soldered onto the substrate, while the fragmented combiner (b) is a single etched copper pattern with no discrete components.  The measured insertion loss (c) confirms that the fragmented combiner achieves lower loss than the COTS version across the entire $L$- to $C$-band operating range.

\begin{figure}
\begin{center}
\includegraphics[angle=0,width=\linewidth]{/Users/jim.maloney/Book/images/ch10/Landgren2018_Fig4_cropped.png}
\caption[Fabricated combiners and measured insertion loss]{Fabricated combiner boards and measured performance: (a) COTS Wilkinson combiner with discrete chips, (b) fragmented combiner with a single etched pattern, and (c) measured insertion loss comparison showing lower loss for the fragmented design across the operating band \cite{Landgren2018Combiner}.}
\label{fig:OES3}
\end{center}
\end{figure}

The fragmented combiner was then integrated with a 4$\times$4 fragmented aperture antenna array.  Landgren et al.\ \cite{Landgren2019Broadband} reported that the complete antenna system---fragmented aperture elements with fragmented combiners and unbalanced feeds---achieved the program performance goals across the 1.4--6.7~GHz bandwidth.  The measured realized gain of the 4$\times$4 array with the fragmented combiner exceeded that of the same array with COTS combiners across the full band.

Figure~\ref{fig:OES4} shows the broadside realized gain of the 4$\times$4 array with the fragmented combiner (blue) compared to the COTS combiner (red).  The fragmented combiner array achieves higher realized gain across the entire $L$- to $C$-band, with the improvement attributable to the lower insertion loss of the fragmented combiner network.  The inset shows the driven layer of the 4$\times$4 antenna array.

\begin{figure}
\begin{center}
\includegraphics[angle=0,width=\linewidth]{/Users/jim.maloney/Book/images/ch10/Landgren2018_Fig5_cropped.png}
\caption[Realized gain: fragmented vs.\ Wilkinson combiner]{Broadside realized gain comparison: fragmented combiner (blue) vs.\ COTS Wilkinson combiner (red) integrated with the same 4$\times$4 fragmented aperture array.  The fragmented combiner achieves higher gain across the operating band due to lower insertion loss.  Inset: driven layer of the antenna array \cite{Landgren2018Combiner}.}
\label{fig:OES4}
\end{center}
\end{figure}

\subsection{Full-Scale Array Validation}

Based on the success of the 4$\times$4 prototype, GTRI fabricated and tested the full 16$\times$16 element array with fragmented combiners.  Landgren et al.\ \cite{Landgren2019Broadband} reported the measured realized gain of the full array, which met the program specifications from $L$-band through $C$-band.

Figure~\ref{fig:OES5} shows the measured realized gain of the full 16$\times$16 array at broadside.  The measured gain (blue) meets or exceeds the program goal (solid black line) across the 1.4--6.7~GHz design bandwidth.  The modeled gain without power combiners (red) represents the theoretical performance of the antenna elements alone; the difference between the red and blue curves is the insertion loss of the combiner network.  The uniform aperture directivity limit (dashed black) provides an upper bound.

\begin{figure}
\begin{center}
\includegraphics[angle=0,width=\linewidth]{/Users/jim.maloney/Book/images/ch10/Landgren2019_Fig5_cropped.png}
\caption[Realized gain of 16$\times$16 array with fragmented combiners]{Measured realized gain of the 16$\times$16 fragmented aperture array with integrated fragmented combiners (blue), compared to the model without combiners (red), the program performance goal (solid black), and the uniform aperture directivity limit (dashed black).  The system meets the program goal across the full 1.4--6.7~GHz bandwidth \cite{Landgren2019Broadband}.}
\label{fig:OES5}
\end{center}
\end{figure}

The fragmented combiner technology demonstrates that the fragmented optimization approach extends naturally from antennas to passive microwave networks.  The same genetic algorithm and FDTD simulation framework that designs the radiating aperture can also design the feed network behind it, producing an integrated system in which every conducting layer---radiator, ground plane, and combiner---is computationally optimized.


\section{Fragmented Metamaterials}
\label{sec:metamaterials}

The fragmented aperture design methodology has also been applied to the design of metamaterials---artificial electromagnetic structures that exhibit properties not found in natural materials.  The fragmented representation used for fragmented aperture antennas is a natural fit for metamaterial design, since metamaterials are typically periodic structures composed of sub-wavelength conducting elements whose geometry determines the effective material properties.

Sun, Jiang, and Wang \cite{Sun2023PixelatedNZIM} applied an elite-preserving genetic algorithm to the design of fragmented near-zero refractive index metamaterials (NZIM).  Near-zero index metamaterials have been proposed for applications including directive emission, cloaking, and tunneling.  The fragmented approach allows the optimizer to discover complex unit cell geometries that achieve the desired near-zero index behavior across a specified frequency band, without requiring the designer to select a specific metamaterial topology (such as split-ring resonators or complementary structures) in advance.  This is the same philosophy as fragmented aperture antenna design: the optimizer discovers the structure, rather than having it imposed by the designer.

Sun et al.\ designed an NZIM unit operating between 8.0 and 8.52~GHz using a 14$\times$14 fragment lattice with quarter-symmetry, reducing the binary chromosome to 28 independent bits.  The unit cell was etched on both sides of a Kappa 438 substrate and optimized using CST Microwave Studio driven by a Python-controlled elite-preserving genetic algorithm (EGA).  Figure~\ref{fig:OES_NZIM1} shows the optimized NZIM unit cell and its extracted electromagnetic properties: the relative permittivity and permeability both approach zero near 8.2~GHz, yielding a refractive index $|n_{\mathrm{re}}| < 1$ across the 8.0--8.52~GHz band.

\begin{figure}
\begin{center}
\includegraphics[angle=0,width=\linewidth]{/Users/jim.maloney/Book/images/ch10/sun2023/Sun2023_Fig4_NZIM_unit_cell.png}
\caption[Fragmented NZIM unit cell and extracted parameters]{Fragmented NZIM design by Sun et al.: (a) optimized 14$\times$14 unit cell on Kappa 438 substrate, (b) extracted relative permittivity and permeability showing near-zero values around 8.2~GHz, and (c) equivalent refractive index demonstrating $|n_{\mathrm{re}}| < 1$ across the 8.0--8.52~GHz band \cite{Sun2023PixelatedNZIM}.}
\label{fig:OES_NZIM1}
\end{center}
\end{figure}

To validate the NZIM performance, Sun et al.\ fabricated a 5$\times$5 array of NZIM unit cells and loaded it as a cladding layer above a rectangular microstrip patch antenna (MPA) resonating at 8.2~GHz.  The near-field simulation results show that the spherical waves radiated by the MPA are transformed into quasi-plane waves after passing through the NZIM cladding, confirming the beam-shaping effect.  Figure~\ref{fig:OES_NZIM2} shows the fabricated prototypes and the near-field distributions with and without the NZIM superstrate.  The field distributions clearly demonstrate the transition from spherical to quasi-planar wavefronts when the NZIM is present.

\begin{figure}
\begin{center}
\includegraphics[angle=0,width=\linewidth]{/Users/jim.maloney/Book/images/ch10/sun2023/Sun2023_Fig5_prototypes_fields.png}
\caption[NZIM prototypes and near-field distributions]{Fabricated MPA and NZIM prototypes with near-field distributions: (a) MPA geometry, (b) side view showing NZIM cladding placement, (c--e) electric field of MPA alone at 8.0, 8.2, and 8.5~GHz showing spherical wavefronts, and (f--h) electric field with NZIM cladding showing quasi-planar wavefronts \cite{Sun2023PixelatedNZIM}.}
\label{fig:OES_NZIM2}
\end{center}
\end{figure}

The measured results confirmed that the NZIM significantly improves the MPA directivity, with an average gain enhancement of 5.3~dB and a maximum gain improvement of 6.05~dB.  Figure~\ref{fig:OES_NZIM3} shows the experimental setup and measured performance: the near-field test system, reflection coefficient comparison, radiation patterns at 5.8~GHz, and realized gain with and without the NZIM cladding.  Compared to previous manually designed NZIMs that achieved 1--5~dB gain enhancement, the automated fragmented design achieves 6.05~dB---a result enabled by the optimizer's ability to explore the full design space of the 14$\times$14 fragment grid.

\begin{figure}
\begin{center}
\includegraphics[angle=0,width=\linewidth]{/Users/jim.maloney/Book/images/ch10/sun2023/Sun2023_Fig6_measurements.png}
\caption[Fragmented NZIM experimental validation and gain enhancement]{Experimental validation of the fragmented NZIM: (a) near-field test system with fabricated prototypes, (b) measured and simulated reflection coefficients, (c) radiation patterns at 5.8~GHz, and (d) realized gain comparison showing up to 6.05~dB gain enhancement with the NZIM cladding \cite{Sun2023PixelatedNZIM}.}
\label{fig:OES_NZIM3}
\end{center}
\end{figure}

As described in Chapter~9, the Periodic Level Set Function (P-LSF) methodology developed by Saad-Falcon et al.\ \cite{SaadFalcon2024LevelSet} has been applied to the design of broadband frequency-selective surfaces and narrowband metasurface filters.  The P-LSF approach is particularly well suited for periodic electromagnetic structures because the distance wrapping technique (Equation~\ref{eq:wrapped} in Chapter~9) ensures that the material distribution is inherently periodic, with no discontinuities at the unit cell boundaries.  The broadband 2:1 bandpass FSS (15--30~GHz) and dual-band high-Q notch filter demonstrated by Saad-Falcon et al.\ represent state-of-the-art examples of computationally optimized metasurface design.

Howard et al.\ \cite{Howard2022LossTangent} applied the P-LSF methodology to a different class of electromagnetic structure: resonant periodic surfaces for dielectric materials characterization.  The optimized surface exhibits a resonance whose magnitude depends linearly on the loss tangent of an adjacent dielectric sample, enabling accurate, non-destructive measurement of dielectric loss.  This is a compelling example of the versatility of the fragmented approach---the same design machinery that produces antennas can also produce precision measurement instruments.

A related and rapidly emerging area is the application of machine learning to digitally programmable metasurfaces, which use fragmented unit cell arrays with electronically controlled phase states to dynamically manipulate electromagnetic waves.  Ayele, Noakoasteen, and Christodoulou \cite{Ayele2025PatternCode} demonstrated the use of CNNs for efficient pattern-to-code mapping of digital metasurfaces: given a desired far-field radiation pattern, the CNN predicts the 2-bit coding sequence for a 15$\times$15 element metasurface operating at 100~GHz.  In a companion study, Ayele et al.\ \cite{Ayele2025SpaceTimeCoding} extended the approach to space-time-coding digital metasurfaces, using a hybrid CNN-LSTM architecture to predict the time-varying coding sequences that produce desired radiation patterns across multiple harmonics.  These works bridge the fragmented metasurface concept with real-time reconfigurability, addressing the computational bottleneck that has limited the practical deployment of programmable metasurfaces for applications including beam steering, RCS reduction, and wireless communications.


\section{Fragmented Absorbers}
\label{sec:absorbers}

Electromagnetic absorbers are structures designed to capture incident electromagnetic energy and convert it to heat, minimizing both reflection and transmission.  Absorbers have important applications in radar cross-section (RCS) reduction, electromagnetic compatibility, and antenna measurement chambers.  The fragmented design approach is well suited to absorber design because the absorption properties depend directly on the geometry and distribution of conducting (and resistive) elements on a surface.

Zhao et al.\ \cite{Zhao2015Novel} demonstrated a novel absorber design based on a fragmented frequency-selective surface (FSS) with integrated lumped resistors, optimized using an estimation of distribution algorithm (EDA).  The absorber consists of four layers: a fragmented FSS layer, a substrate, a dielectric (honeycomb) spacer, and a metallic ground plane.  The FSS unit cell is fragmented into a 10$\times$10 grid of sub-wavelength conducting patches, with lumped resistors placed at selected locations between adjacent fragments.  Figure~\ref{fig:OES_absorber1} shows the absorber layer structure and the simulated unit cell geometry.

\begin{figure}
\begin{center}
\includegraphics[angle=0,width=0.75\linewidth]{/Users/jim.maloney/Book/images/ch10/zhao2015/Zhao2015_Fig1_2_structure_unitcell.png}
\caption[Fragmented FSS absorber structure and unit cell]{Fragmented FSS absorber by Zhao et al.: (a) side view showing the four-layer structure (FSS, substrate, dielectric spacer, ground plane), and (b) top view of the 10$\times$10 fragmented unit cell with conducting patches and resistor locations \cite{Zhao2015Novel}.}
\label{fig:OES_absorber1}
\end{center}
\end{figure}

The key innovation of Zhao et al.'s approach is the use of EDA rather than a conventional genetic algorithm.  EDA builds a probabilistic model of the design space and samples new candidates from the learned distribution, which is particularly effective for problems where the positions of metallic fragments and resistors are correlated.  Three parameters are jointly optimized: the metallic fragment positions ($\mathbf{F}_{\mathrm{position}}$), the resistor positions ($\mathbf{R}_{\mathrm{position}}$), and the resistor values ($\mathbf{R}_{\mathrm{value}}$).  The fitness function balances the absorption bandwidth (below $-6$~dB reflectivity) with the mean reflectivity across the 2--6~GHz target band.

Figure~\ref{fig:OES_absorber2} shows the optimized unit cell geometries and reflectivity performance.  The fragmented FSS with resistors (a) achieves two strong absorption peaks at 3.64~GHz ($-8.67$~dB) and 5.86~GHz ($-9.77$~dB), while the version without resistors (b) exhibits only weak absorption.  The loaded resistors fundamentally change the absorbing mechanism by providing the necessary energy dissipation.

\begin{figure}
\begin{center}
\includegraphics[angle=0,width=\linewidth]{/Users/jim.maloney/Book/images/ch10/zhao2015/Zhao2015_Fig5_optimized_reflectivity.png}
\caption[Optimized FSS unit cells and simulated reflectivity]{Optimized fragmented FSS unit cells and performance: (a) with resistors, (b) without resistors, and (c) corresponding simulated reflectivity curves showing two strong absorption peaks for the resistor-loaded design.  Insets show surface current distributions at the two absorption peaks \cite{Zhao2015Novel}.}
\label{fig:OES_absorber2}
\end{center}
\end{figure}

The optimized absorber was fabricated and measured in a microwave anechoic chamber.  Figure~\ref{fig:OES_absorber3} shows the fabricated samples and measured reflectivity curves.  The measured absorption band ($-6$~dB reflectivity) extends from 3.13 to 6.42~GHz, with reflectivity values of $-14.39$~dB at 3.80~GHz and $-10.14$~dB at 5.68~GHz.  The total thickness of the absorber is only 4~mm, corresponding to approximately $0.04\lambda$ at the lowest operating frequency and $0.08\lambda$ at the highest---making this an ultra-thin wideband absorber.  The measured results fit well with simulations, validating the EDA-based fragmented design approach.

\begin{figure}
\begin{center}
\includegraphics[angle=0,width=\linewidth]{/Users/jim.maloney/Book/images/ch10/zhao2015/Zhao2015_Fig6_fabricated_measured.png}
\caption[Fabricated FSS absorber and measured reflectivity]{Fabricated fragmented FSS absorber samples: (a) with resistors, (b) without resistors, and (c) measured and simulated reflectivity curves confirming wideband absorption from 3.13 to 6.42~GHz \cite{Zhao2015Novel}.}
\label{fig:OES_absorber3}
\end{center}
\end{figure}

Lee et al.\ \cite{Lee2023UWBAbsorber} extended the fragmented absorber concept to ultra-wideband operation by using two stacked GA-optimized metasurface layers separated by PMI foam dielectric spacers, backed by a carbon fiber reinforced polymer (CFRP) reflector.  Each metasurface layer uses a 20$\times$20 fragment grid with x/y-axis symmetry (55 independent fragments per layer), printed with carbon paste via silk-screen printing on polyimide film.  The dual-layer approach creates multiple destructive interference paths that broaden the absorption bandwidth dramatically.

Figure~\ref{fig:OES_absorber4} shows the optimized metasurface patterns and simulated reflectance performance.  The GA-optimized design achieves greater than 90\% absorption (reflectance below $-10$~dB) from 2.2 to 18~GHz---a 156\% fractional bandwidth covering the entire S- through Ku-band range.

\begin{figure}
\begin{center}
\includegraphics[angle=0,width=\linewidth]{/Users/jim.maloney/Book/images/ch10/lee2023/Lee2023_Fig4_patterns_reflectance.png}
\caption[Dual-layer fragmented absorber patterns and reflectance]{GA-optimized dual-layer fragmented absorber by Lee et al.: (a) upper and lower metasurface fragment patterns (20$\times$20 grids with x/y symmetry), and (b) simulated reflectance vs.\ sheet resistance, showing $>$90\% absorption from 2.2 to 18~GHz at the optimal 40~$\Omega$/sq \cite{Lee2023UWBAbsorber}.}
\label{fig:OES_absorber4}
\end{center}
\end{figure}

The absorber was fabricated as a composite structure with a glass fiber reinforced polymer (GFRP) protective outer layer, making it suitable for aerospace applications.  Figure~\ref{fig:OES_absorber5} shows the fabricated metasurface layers and the assembled composite absorber.  The prototype passed MIL-STD-810H environmental durability testing (temperature cycling and humidity exposure), demonstrating the practical viability of fragmented absorber technology for real-world deployment.  The total thickness is 13~mm ($0.06\lambda_0$ at 2~GHz), and absorption remains above 90\% for incidence angles up to 30$^\circ$ in both TE and TM polarizations.

\begin{figure}
\begin{center}
\includegraphics[angle=0,width=\linewidth]{/Users/jim.maloney/Book/images/ch10/lee2023/Lee2023_Fig7_fabricated.png}
\caption[Fabricated fragmented composite absorber assembly]{Fabricated fragmented composite absorber: (a) silk-screen printed upper and lower metasurface patterns on polyimide film, and (b) assembled composite absorber showing top view and side view with layer stack (GFRP, metasurfaces, PMI foam, CFRP reflector) \cite{Lee2023UWBAbsorber}.}
\label{fig:OES_absorber5}
\end{center}
\end{figure}


\section{Fragmented Waveguide Components}
\label{sec:waveguide}

The fragmented design concept has also been applied to waveguide-based structures, particularly substrate integrated waveguide (SIW) horn antennas.  SIW technology implements conventional rectangular waveguide structures in planar printed circuit board form using rows of metallized vias as the waveguide sidewalls, enabling low-cost fabrication and easy integration with planar circuits.  However, SIW horn antennas typically suffer from narrow bandwidth, high sidelobe levels, and unequal beamwidths in the E- and H-planes.

Zechmeister, Lacik, and Kadlec \cite{Zechmeister2021Design} applied fragmented optimization to the design of SIW horn antennas operating in the 24~GHz ISM band.  The antenna aperture is extended beyond the conventional flared horn and overlaid with a fragmented grid of dielectric and air-filled cells.  By optimizing the binary pattern of filled and empty cells, the radiation properties---beamwidth, sidelobe level, gain, and impedance match---can be controlled simultaneously.

The optimization used a novel heuristic algorithm called Binary Ink Stamp Optimization (BISO), which is inspired by stamps with different ink saturation levels.  BISO reflects the significance of individual bits during the optimization process, giving it an advantage over standard genetic algorithms for binary pattern problems.  The fitness function combines terms for reflection coefficient, beamwidth equality in the E- and H-planes, sidelobe level, and gain.  Figure~\ref{fig:OES_SIW1} shows the simplified model of the fragmented SIW horn antenna and the optimized design with the SIW feeding structure.

\begin{figure}
\begin{center}
\includegraphics[angle=0,width=\linewidth]{/Users/jim.maloney/Book/images/ch10/zechmeister2021/Zech2021_Fig4_6_LP_design.png}
\caption[Fragmented SIW horn antenna design]{Fragmented SIW horn antenna design by Zechmeister et al.: (left) simplified model showing the fragmented region applied over the horn aperture, and (right) optimized linearly polarized antenna with SIW structure and coaxial-SIW transition.  The dielectric part is 3D-printed and the metal layer is cut copper foil \cite{Zechmeister2021Design}.}
\label{fig:OES_SIW1}
\end{center}
\end{figure}

Two antenna designs were demonstrated: a linearly polarized (LP) version and a circularly polarized (CP) version using an antipodal configuration.  Both use a single-layer structure---a key advantage over multi-layered approaches---and are fabricated using low-cost 3D printing for the dielectric substrate with copper foil for the metallization.

The LP antenna achieves a measured gain of 14.2~dBi at the 24~GHz ISM band center frequency, with reflection coefficient below $-18$~dB across the full band.  The E- and H-plane half-power beamwidths are 26.5$^\circ$ and 25$^\circ$, respectively---nearly equal, which is the primary design objective.  Figure~\ref{fig:OES_SIW2} shows the manufactured prototype, and Figure~\ref{fig:OES_SIW3} shows the measured radiation patterns confirming good agreement with simulation and nearly equal beamwidths.

\begin{figure}
\begin{center}
\includegraphics[angle=0,width=\linewidth]{/Users/jim.maloney/Book/images/ch10/zechmeister2021/Zech2021_Fig8_LP_prototype.png}
\caption[LP fragmented SIW horn prototype]{Manufactured prototype of the linearly polarized fragmented SIW horn antenna.  The dielectric is 3D-printed XT co-polyester and the metal is cut copper foil with soldered vias \cite{Zechmeister2021Design}.}
\label{fig:OES_SIW2}
\end{center}
\end{figure}

\begin{figure}
\begin{center}
\includegraphics[angle=0,width=\linewidth]{/Users/jim.maloney/Book/images/ch10/zechmeister2021/Zech2021_Fig10_LP_patterns.png}
\caption[LP fragmented SIW horn radiation patterns]{Measured and simulated radiation patterns of the linearly polarized fragmented SIW horn antenna at 24.125~GHz: (a) E-plane and (b) H-plane.  The nearly equal half-power beamwidths (26.5$^\circ$ and 25$^\circ$) confirm the effectiveness of the fragmented optimization \cite{Zechmeister2021Design}.}
\label{fig:OES_SIW3}
\end{center}
\end{figure}

The CP version uses an antipodal structure with top and bottom metal layers to generate circular polarization.  The fragment grid uses 6$\times$28 cells with 13 parameterized column lengths to control the fragment pattern shape.  The measured CP antenna achieves 8.9~dBi RHCP gain, an axial ratio below 3~dB across the 24~GHz ISM band, sidelobe levels below $-7.4$~dB, and half-power beamwidths of 30.6$^\circ$ and 33.4$^\circ$ in the two principal planes.  Figure~\ref{fig:OES_SIW4} shows the manufactured CP prototype.

\begin{figure}
\begin{center}
\includegraphics[angle=0,width=\linewidth]{/Users/jim.maloney/Book/images/ch10/zechmeister2021/Zech2021_Fig15_CP_prototype.png}
\caption[CP fragmented SIW horn prototype]{Manufactured prototype of the circularly polarized fragmented SIW horn antenna, shown alongside a coin for scale.  The antipodal structure achieves RHCP with 8.9~dBi gain and axial ratio below 3~dB across the 24~GHz ISM band \cite{Zechmeister2021Design}.}
\label{fig:OES_SIW4}
\end{center}
\end{figure}

In a follow-up study, Zechmeister and Lacik \cite{Zechmeister2022Pixelated} improved the CP design by replacing the binary fragmentation with a four-state scheme (air, 0.6~mm, 1.0~mm, and 1.6~mm diameter holes), providing finer control over the dielectric distribution.  The improved design is 29\% shorter than the original, with a wider $-10$~dB reflection coefficient bandwidth of 21.4\%, a 3~dB axial ratio bandwidth of 4.6\%, and improved beamwidth equality (less than 1$^\circ$ difference between principal planes).

The fragmented SIW horn work demonstrates that the fragmented aperture philosophy extends naturally to volumetric waveguide structures.  Rather than optimizing a planar conducting pattern, the optimizer controls the three-dimensional distribution of dielectric material within the horn, achieving performance that would be difficult to obtain through conventional parametric design.


\section{Fragmented Radomes and RCS Reduction}
\label{sec:radomes}


A particularly compelling application of fragmented and optimized metasurface design is in the development of low-observable radomes---protective enclosures for antennas that must transmit in-band signals efficiently while minimizing out-of-band radar cross section (RCS).  Conventional frequency-selective surface (FSS) radomes reflect out-of-band energy back toward the source, creating a significant RCS signature.  By combining fragmented phase gradient metasurfaces with bandpass FSS structures, and optimizing the composite design with algorithms such as particle swarm optimization (PSO), it is possible to redirect out-of-band energy away from the monostatic direction while maintaining high in-band transmission.

Haji-Ahmadi et al.\ \cite{HajiAhmadi2017Checkerboard} demonstrated a fragmented checkerboard metasurface for ultra-wideband RCS reduction.  The design uses a chessboard arrangement of two types of fragmented artificial magnetic conductor (AMC) unit cells, each optimized via binary particle swarm optimization (BPSO) to produce a specific reflection phase.  Each cell uses a 16$\times$16 fragment grid with 8-fold symmetry, reducing the optimization to 36 independent bits per cell type (72 total).  The two cell types are arranged in a checkerboard pattern such that their reflected fields maintain a 180$^\circ \pm 37^\circ$ phase difference, causing destructive interference in the broadside direction and scattering the energy into other angles.

Figure~\ref{fig:OES_RCS1} shows the optimized fragmented unit cell layouts and the resulting out-of-phase reflection bandwidth.  The BPSO optimizer discovers irregular, non-intuitive conducting patterns that would be impossible to design analytically.

\begin{figure}
\begin{center}
\includegraphics[angle=0,width=\linewidth]{/Users/jim.maloney/Book/images/ch10/hajiahmadi2017/HajiAhmadi2017_Fig7_cells_phase.png}
\caption[Fragmented checkerboard RCS reducer unit cells]{Fragmented checkerboard RCS reducer by Haji-Ahmadi et al.: (a) optimized fragmented unit cell layouts for cell type 1 and cell type 2, and (b) reflection phase from each cell and phase difference, showing the required 180$^\circ \pm 37^\circ$ out-of-phase condition maintained from 3.6 to 10.4~GHz \cite{HajiAhmadi2017Checkerboard}.}
\label{fig:OES_RCS1}
\end{center}
\end{figure}

The fragmented checkerboard metasurface was fabricated on Rogers RO4003c substrate with a PTFE spacer and measured in an anechoic chamber.  Figure~\ref{fig:OES_RCS2} shows the fabricated prototype and the measured monostatic RCS comparison.  The measured 10~dB RCS reduction bandwidth extends from 3.8 to 10.7~GHz (95\% fractional bandwidth), with a maximum reduction of 36.5~dB at 7.3~GHz.  This significantly outperforms prior checkerboard designs based on conventional (non-fragmented) AMC unit cells, which achieved fractional bandwidths of only 41--63\%.

\begin{figure}
\begin{center}
\includegraphics[angle=0,width=\linewidth]{/Users/jim.maloney/Book/images/ch10/hajiahmadi2017/HajiAhmadi2017_Fig13_RCS_measured.png}
\caption[Measured monostatic RCS of checkerboard metasurface]{Measured and simulated monostatic RCS of the fragmented checkerboard metasurface, normalized to an equal-size metal plate.  The measured 10~dB RCS reduction bandwidth extends from 3.8 to 10.7~GHz (95\% fractional bandwidth), with a maximum reduction of 36.5~dB at 7.3~GHz \cite{HajiAhmadi2017Checkerboard}.}
\label{fig:OES_RCS2}
\end{center}
\end{figure}

Pang et al.\ \cite{Pang2024Radome} proposed a broadband low-scattering radome that combines phase gradient metasurfaces (PGMs) with a frequency-selective surface (FSS).  The key innovation is a single-layer composite structure: a modified Jerusalem cross resonator provides a phase gradient across six unit cells (each contributing $60^\circ$ of phase shift), while a non-uniform bandpass FSS replaces the conventional metallic ground plane to allow in-band signal transmission.  The PGM redirects out-of-band reflected energy away from the monostatic direction via the generalized Snell's law, while the FSS acts as a high-pass filter transmitting in-band signals.  Particle swarm optimization (PSO) was applied to jointly optimize the FSS slot dimensions for maximum RCS reduction bandwidth while preserving high in-band transmission.

Figure~\ref{fig:OES_PangStructure} shows the PGMs-FSS composite structure.  The six unit cells, each with different dimensions of the Jerusalem cross resonator, are arranged in a one-dimensional gradient that deflects the reflected wave away from the specular direction.

\begin{figure}
\begin{center}
\includegraphics[angle=0,width=\linewidth]{/Users/jim.maloney/Book/images/ch10/pang2024/Pang2024_Fig9_PGMs_FSS_structure.png}
\caption[PGMs-FSS radome structure]{PGMs-FSS radome structure by Pang et al.: (a) overall composite structure, (b) phase gradient metasurface layer, (c) non-uniform FSS layer with six different slot geometries, and (d) comparison of RCS reduction for the non-uniform FSS alone versus the combined PGMs-FSS \cite{Pang2024Radome}.}
\label{fig:OES_PangStructure}
\end{center}
\end{figure}

A prototype measuring $288 \times 288 \times 3$~mm$^3$ (thickness $0.08\lambda_l$, where $\lambda_l$ is the wavelength at the lowest operating frequency) was fabricated and measured.  Figure~\ref{fig:OES_PangResults} shows the fabricated prototype and the measured versus simulated RCS reduction.  The measured results demonstrate over 10~dB RCS reduction from 7.95 to 18.12~GHz, corresponding to a 78\% relative bandwidth spanning the X and Ku bands.  Insertion loss remains below 1~dB within the 15.95--17.7~GHz passband.  This performance significantly exceeds that of conventional single-layer FSS radomes, which typically achieve narrower RCS reduction bandwidths.

\begin{figure}
\begin{center}
\includegraphics[angle=0,width=\linewidth]{/Users/jim.maloney/Book/images/ch10/pang2024/Pang2024_Fig14_prototype.png}
\caption[Fabricated PGMs-FSS radome prototype]{(a) Fabricated PGMs-FSS radome prototype ($288 \times 288 \times 3$~mm$^3$) and (b) measurement setup in anechoic chamber using standard horn antennas \cite{Pang2024Radome}.}
\label{fig:OES_PangPrototype}
\end{center}
\end{figure}

\begin{figure}
\begin{center}
\includegraphics[angle=0,width=0.85\linewidth]{/Users/jim.maloney/Book/images/ch10/pang2024/Pang2024_Fig15_RCS_measured.png}
\caption[Measured RCS reduction of PGMs-FSS radome]{Measured and simulated RCS reduction of the PGMs-FSS radome under normal incidence.  Over 10~dB RCS reduction is achieved from 7.95 to 18.12~GHz (78\% relative bandwidth), with good agreement between measurement and simulation \cite{Pang2024Radome}.}
\label{fig:OES_PangResults}
\end{center}
\end{figure}

More recently, Zeng et al.\ \cite{Zeng2026ChessboardRCS} applied a CNN-KAN (convolutional neural network--Kolmogorov-Arnold network) deep learning framework to accelerate inverse design of chessboard metasurfaces for RCS reduction.  By training on a dataset of fragmented unit cell geometries and their full-wave simulated reflection phases, the surrogate model replaces iterative electromagnetic simulation during optimization, reducing design time from hours to seconds.  The resulting optimized checkerboard metasurface achieved greater than 10~dB monostatic RCS reduction from 7.4 to 15.2~GHz (69\% relative bandwidth), demonstrating that machine learning can effectively navigate the vast design space of fragmented surfaces.

These radome and RCS reduction designs illustrate how the fragmented optimization approach addresses a fundamentally different electromagnetic challenge: rather than maximizing radiation or absorption, the goal is to simultaneously control transmission (in-band) and scattering (out-of-band) properties of a surface.


\section{Fragmented Energy Harvesting}
\label{sec:harvesting}

The fragmented metasurface concept has also been applied to electromagnetic energy harvesting---the capture of ambient or intentionally transmitted RF energy and its conversion to usable DC power.  Energy harvesting metasurfaces must efficiently absorb incident electromagnetic energy across multiple frequency bands and polarizations, then couple it to resistive loads (or rectifying circuits) through vias.

Ghaderi et al.\ \cite{Ghaderi2018EnergyHarvesting} demonstrated a fragmented metasurface for dual-band, multi-polarization electromagnetic energy harvesting.  The unit cell combines a ring resonator (for the 2.45~GHz WiFi band) and a symmetric electric-LC (ELC) resonator (for the 6~GHz band), with resistive loads connected through vias to the ground plane.  The cell surface is fragmented into a 27$\times$27 grid (0.41~mm fragment resolution) with 4-fold symmetry, yielding approximately 105 independent optimization variables.  The BPSO algorithm simultaneously optimizes the fragment pattern, load resistance ($R = 240~\Omega$), and gap size ($g = 0.25$~mm).

Figure~\ref{fig:OES_harvest1} shows the BPSO-optimized fragmented unit cell.  The irregular, fractal-like conducting pattern discovered by the optimizer combines the dual-band resonant behavior with efficient coupling to the via-connected loads.  The design achieves polarization-independent operation---a critical requirement for harvesting ambient RF energy arriving from arbitrary directions and polarizations.

\begin{figure}
\begin{center}
\includegraphics[angle=0,width=\linewidth]{/Users/jim.maloney/Book/images/ch10/ghaderi2018/Ghaderi2018_Fig5_6_pixelated_cell.png}
\caption[Fragmented energy harvesting unit cell]{BPSO-optimized fragmented energy harvesting unit cell by Ghaderi et al.: (top) 3D view of the optimized 27$\times$27 fragment pattern showing the irregular conducting geometry with vias to loads and edge capacitors, and (bottom) top and bottom views of the fabrication-ready design with 240~$\Omega$ resistive loads \cite{Ghaderi2018EnergyHarvesting}.}
\label{fig:OES_harvest1}
\end{center}
\end{figure}

A 9$\times$9 array of the optimized unit cells was fabricated on Rogers RT/duroid 6006 substrate and measured in an anechoic chamber.  Figure~\ref{fig:OES_harvest2} shows the measurement setup.  The measured harvesting efficiency exceeds 90\% at both 2.45 and 6~GHz for all polarization angles, confirming the polarization-independent design.  The delivered power is also insensitive to incidence angle from 0$^\circ$ to 45$^\circ$ for both TE and TM polarizations.  This work demonstrates that the fragmented optimization approach can produce electromagnetic structures optimized for energy capture rather than radiation or scattering---yet another application of the same fundamental design philosophy.

Keshavarz et al.\ \cite{Keshavarz2024PixelatedRectenna} applied fragmented antenna design to rectenna (rectifying antenna) systems for wireless power transfer.  In a conventional rectenna, a matching network is required between the receiving antenna and the rectifier diode to conjugate-match their impedances; this network adds loss, size, and design complexity.  By optimizing the fragment pattern of the receiving antenna using binary particle swarm optimization (BPSO) with a multi-objective fitness function that simultaneously targets resonant frequency, conjugate impedance match to the diode, and antenna gain, the authors eliminated the matching network entirely.  The BPSO discovers a non-intuitive fragment pattern that provides the correct complex impedance at the antenna terminals to directly drive the rectifier.  The resulting fragmented rectenna achieved 37\% RF-to-DC conversion efficiency at 0~dBm input power (2.5~GHz), rising to 64\% at $+12$~dBm, in a compact single-layer structure---demonstrating that the fragmented aperture approach can co-optimize antenna radiation and circuit-level impedance matching in a single design step.

\begin{figure}
\begin{center}
\includegraphics[angle=0,width=\linewidth]{/Users/jim.maloney/Book/images/ch10/ghaderi2018/Ghaderi2018_Fig14_measurement.png}
\caption[Energy harvesting metasurface measurement setup]{Measurement of the fragmented energy harvesting metasurface: (a) schematic of the test setup with 50-to-240~$\Omega$ converter and spectrum analyzer, and (b) anechoic chamber photograph showing the 9$\times$9 array under test with a transmitter horn antenna \cite{Ghaderi2018EnergyHarvesting}.}
\label{fig:OES_harvest2}
\end{center}
\end{figure}


\section{Fragmented Microwave Filters}
\label{sec:filters}

Microwave filters are fundamental components in communication systems, and their design has traditionally relied on analytical synthesis methods based on coupled resonator theory.  Gomez et al.\ \cite{Gomez2025PixelatedFilters} demonstrated that the fragmented inverse design approach can produce compact bandpass and notch filters in coplanar waveguide (CPW) technology, achieving performance comparable to conventional designs in a significantly smaller footprint.

The filter design region is a fragmented area within a CPW transmission line, where each fragment is either conducting (copper) or absent (etched away).  The optimization uses a genetic algorithm combined with direct binary search and randomization mutation, driven by CST Microwave Studio simulations.  The design begins with a coarse 24$\times$16 grid (400~$\mu$m fragments) and refines to a 48$\times$32 grid (200~$\mu$m fragments).  A key innovation is the integration of short-circuited stubs at the top and bottom of the fragmented region, which dramatically improves stopband rejection.

Figure~\ref{fig:OES_filter1} shows the proposed fragmented notch filter schematic and the fabricated prototype.  The filter is fabricated on Rogers Kappa 438 substrate via UV laser ablation, which precisely removes copper to define the fragmented pattern at 200~$\mu$m resolution.  The irregular conducting pattern discovered by the optimizer bears no resemblance to conventional coupled-line or interdigitated filter topologies.

\begin{figure}
\begin{center}
\includegraphics[angle=0,width=\linewidth]{/Users/jim.maloney/Book/images/ch10/gomez2025/Gomez2025_Fig1_filter_photo.png}
\caption[Fragmented microwave notch filter design and prototype]{Fragmented microwave notch filter by Gomez et al.: (a) schematic showing the irregular optimized conducting pattern in a CPW structure, and (b) photograph of the fabricated filter on Rogers Kappa 438 substrate with SMA connectors \cite{Gomez2025PixelatedFilters}.}
\label{fig:OES_filter1}
\end{center}
\end{figure}

Figure~\ref{fig:OES_filter2} shows the measured and simulated S-parameters.  The filter without short-circuited stubs achieves a notch depth of $-26$~dB at 5~GHz, while the version with shorted stubs reaches $-61$~dB---a 35~dB improvement.  The fragmented filter occupies only 61.5~mm$^2$, less than half the area of an equivalent third-order Chebyshev interdigitated filter.  This compact size, combined with the single-layer, via-free construction, makes fragmented filters attractive for integrated microwave circuits.

Karahan, Liu, and Sengupta \cite{Karahan2023InverseMMWave} demonstrated that the fragmented inverse design concept extends to on-chip millimeter-wave circuits.  By discretizing a 300$\times$300~$\mu$m area into a 16$\times$16 fragment grid and training a deep CNN forward model on the scattering parameters, they enabled rapid synthesis of multi-port EM structures with designer S-parameters.  The approach was applied to the co-design of a broadband 30--100~GHz output matching network with a SiGe power amplifier, achieving 16\%--24.7\% power-added efficiency across a 103.2\% fractional bandwidth---performance that would be extremely difficult to achieve through conventional template-based circuit design.  This work demonstrates that the fragmented optimization philosophy, originally developed for antenna apertures, can produce globally competitive designs in the domain of integrated circuits.

\begin{figure}
\begin{center}
\includegraphics[angle=0,width=\linewidth]{/Users/jim.maloney/Book/images/ch10/gomez2025/Gomez2025_Fig7_Sparams.png}
\caption[Fragmented filter S-parameters and notch depth]{Measured and simulated S-parameters of the fragmented filter: (top) without short-circuited stubs showing $-26$~dB notch depth, and (bottom) with shorted stubs achieving $-61$~dB notch depth---a 35~dB improvement.  The fragmented design area is less than half the size of an equivalent conventional filter \cite{Gomez2025PixelatedFilters}.}
\label{fig:OES_filter2}
\end{center}
\end{figure}


\section{Photonic Integration with Fragmented Arrays}
\label{sec:photonic}

An emerging application of fragmented aperture antennas is their integration with high-power photodiodes for fiber-fed RF arrays.  In conventional phased array systems, RF signals are distributed to antenna elements through coaxial cables or waveguides, which become increasingly lossy and heavy at higher frequencies and over longer runs.  Replacing these metallic feed lines with optical fiber offers dramatic reductions in size, weight, and power (SWaP): fiber is lighter, lower-loss over long distances, and immune to electromagnetic interference.  The challenge is integrating the photodiode---which converts the modulated optical signal back to RF---directly at the antenna element without degrading the antenna's wideband performance.

\begin{figure}
\begin{center}
\includegraphics[angle=0,width=\linewidth]{/Users/jim.maloney/Book/images/ch10/hughes2015/Hughes2015_Fig1_XBandArray.png}
\caption[X-band fragmented array with photodiode integration]{X-band fragmented aperture array integrated with high-power photodiodes.  (a) Design goals and measured performance (gain and port reflection vs.\ frequency); (b) rendering of a unit cell showing the fragmented conducting layer, substrate, and feed; (c) fabricated 5$\times$5 array on Rogers PCB; (d) close-up of the fragmented aperture pattern; (e) measured horizontal scan pattern showing $\sim\!13$~dB sidelobe suppression across 5--14~GHz \cite{Hughes2015Integration}.}
\label{fig:OES_photonic1}
\end{center}
\end{figure}

Hughes et al.\ at GTRI and the University of Virginia demonstrated this integration in two complementary efforts \cite{Hughes2015Co-design,Hughes2015Integration}.  The fragmented aperture array was designed for linear polarization, 50-$\Omega$ input impedance, $>$10~dB return loss, and flat gain response across 6--12~GHz (Figure~\ref{fig:OES_photonic1}).  The GA-optimized unit cell was tessellated into a 5$\times$5 array on low-cost Rogers PCB.  At 12~GHz, the array demonstrated approximately 10~dBi gain with 20$^\circ$ beamwidth, and measured gain mismatch between any two elements was less than 5~dB across 5--14~GHz.

The UVA-developed InGaAs/InP modified uni-traveling carrier (MUTC) photodiodes exhibited $\sim\!10$~GHz RF bandwidth and CW output power exceeding 1~W---sufficient to enable net positive RF gain over the photonic link.  GTRI developed a custom coupon board packaging approach that co-locates four active photodiodes with the antenna elements, using fiber-attached lenses for optical injection through a quartz window.  The integrated photodiode-antenna assembly demonstrated passband operation from 8.2 to 12.4~GHz with less than 6~dB gain variation.

The fragmented aperture design proved particularly well-suited to photonic integration because the optimizer could account for the photodiode packaging in the electromagnetic model, ensuring that the nearby photonic structures did not degrade antenna performance.  This flexibility---the ability to co-optimize the antenna fragment pattern in the presence of non-traditional feed structures---is a distinctive advantage of the fragmented aperture approach over conventional antenna topologies.


\section{Summary}

The examples in this chapter demonstrate that the fragmented aperture concept is not limited to antenna design.  The core methodology---fragmented conducting surfaces optimized by evolutionary algorithms with full-wave electromagnetic evaluation---applies to any electromagnetic structure whose performance depends on the spatial distribution of conducting material:

\begin{itemize}
\item \textbf{Broadband phase surfaces} with engineered transmission phase properties, as demonstrated in the original patent.
\item \textbf{Reflectarrays} with similarity-optimized fragmented elements achieving 97\% aperture phase efficiency and the ability to simultaneously display optical images.
\item \textbf{RF power combiners} that outperform conventional COTS components in insertion loss and bandwidth, enabling integrated all-fragmented antenna systems.
\item \textbf{Fragmented absorbers} that achieve ultra-thin, wideband absorption through jointly optimized conducting fragments and lumped resistors.
\item \textbf{Fragmented waveguide components}, specifically SIW horn antennas with optimized dielectric distributions that achieve equal beamwidths and circular polarization in single-layer structures.
\item \textbf{Radomes and RCS reduction} surfaces that combine optimized metasurfaces with frequency-selective structures for simultaneous in-band transmission and out-of-band scattering control.
\item \textbf{Energy harvesting} metasurfaces and rectennas, including fragmented designs that eliminate matching networks by co-optimizing antenna impedance with rectifier diode requirements.
\item \textbf{Microwave filters} with fragmented conducting patterns that achieve compact notch and bandpass responses in single-layer CPW structures.
\item \textbf{Metamaterials and metasurfaces} including near-zero index materials with measured 6~dB gain enhancement, broadband frequency-selective surfaces, and resonant structures for materials characterization.
\item \textbf{Photonic integration} of fragmented aperture arrays with high-power photodiodes for fiber-fed RF systems, exploiting the design flexibility of the GA-optimized fragment pattern to accommodate non-traditional feed structures.
\end{itemize}

As the fragmented aperture design process continues to be refined---through improved optimization algorithms (Chapter~3), continuous parameterizations (Chapter~9), and machine learning acceleration (Chapter~9)---its application to non-antenna electromagnetic structures will expand further.  Any electromagnetic problem that can be cast as ``what pattern of conducting material on a surface optimizes a given performance metric?'' is a candidate for the fragmented aperture approach.


\FloatBarrier

% Bibliography for this chapter
\bibliography{../Literature/master_bibliography,%
              ../Literature/fragmented_aperture_core,%
              ../Literature/ml_ai_optimization_methods,%
              ../Literature/specialized_applications,%
              ../Literature/recent_innovations,%
              ../Literature/metamaterials_ch7}
\bibliographystyle{IEEEtran}
